\chapter{Introduction}
\label{chapter:introduction}


\section{Motivation}

Design thesis, rather than research thesis. General objective: design and test a workflow that answers to a specific need + approach requested by a research group for which this project is done. Partners and participation. Broader context on the fact that this is an aerospace engineering project at the service of the medical field.



* Questo è l'obiettivo di questa tesi: fare il prossimo passo, cioè avviare il progetto concreto basato sulle analisi di concetto fatte fino a questo momento. La tesi si concentra dunque su un singolo looop progettuale all'interno del loop di progetto aeronautico. Il focus è la progettazione e la simulazione di un dispositivo che protegga il payload in caso di failure di drone e sgancio dallo stesso. Il dispositivo ha come obiettivo quello di rendere le accelerazioni di impatto date dalla missione in caso di dropping con gli ammissibili del contenitore e delle fialette mediche trasportate, iin rispetto delle normative che si preoccupano del danno causato dall'impatto su zone civili e dispersione delle sostanze chimiche nell'ambiente. 

Il vantaggio di questo approccio, cioè sul partire dal dispositivo di safety, è l'impostazione del progetto modulare. Questo dispositivo, se realizzato con certi criteri di safety, è innestabile su droni adatti alle missioni pensate in un'ottica modulare, in cui una volta progettato il dispositivo di protezione, si va a convergenza di progetto modificando i requisiti di missione. 

\section{Project framing}
\subsection{The need for improved logistics in the Italian NHS}
logistic failures and why people have thought of drone networks, with some mentions to UAV applications in the medical field and a focused Italian market analysis
* scenario di una sanità devastata, inefficiente, con problemi di logistica e soddisfazione scarsa delle persone, con problema evidenziato ancora di più da crisi sanitaria come il covid, che nonostante sia avvenuto nel 2020-2021 ha visto Torino come uno dei protagonisti principali, prostrando gli ospedali e la popolazione. 

\subsection{PoliTo’s research group’s work (including normative) and expected contribution}

* Iniziativa dell'A.o.U. città della salute e della scienza di Torino, in collaborazione con il Politecnico di Torino e in occasione del presente lavoro di tesi anche con FEV Italia S.r.l, ha preso l'iniziativa di ricerca per tentare di risolvere il problema, iniziando una serie di progetti a catena con una visione a lungo termine: snellire i l trasporto logistico all'interno della sanità implementando un sistema di droni, che possano aiutare con vari scenari. 

Knowing that the long-term objective is the drone network, this is the work that this particular research group has done until now. The 5 MSc theses, the scientific works, etc. My contribution: an answer to PoliTo’s research: formalise (at a high level, non-detailed level) the project and take a little step forward, toward the long term objective, deciding on the approach. 


* Questo progetto vede l'intersezione tra l'aeronautica e la medicina. Mettere a disposizione una tecnologia aerea per contribuire al miglioramento del sistema sanitario nazionale italiano. 
* Sin da subito, il progetto ha assunto un focus logistico, come ogni progetto aeronautico. Impostandolo su analisi di sicurezza, analisi di mercato su scelta di drone, analisi di requisiti, etc. (fai RIASSUNTO) scelte di percoros ottimale, analisi di missione, etc. Questo ha fatto sì che si cominciasse a delineare uno scenario in cui è possibile ottenere l'approvazione dell'ENAC, ente nazionale bla bla bla, per formalizzare questa tipologia di operazione. D'altro canto, in uno scenario in cui per realizzare un'operazione del genere serve la certificazione aeronautica per evidenti rischi causati dal trasporto di medicinali su zone intensamente abitate come Torino, è necessario cominciare una progettazione più concreta per avviare un design loop guidato dalle stringenti normative. 
